% Pro-category completion, developed during the independent configuration audit.
% Provisional B.PRO identifiers; parent editor may renumber.
\subsection{Completion of ordinary pro-reconstruction in arbitrary transcendence degree}

We retain the ordinary category \(\mathsf T\), the model pro-objects,
the slice over \(\mathcal B\), and the limit-point definition of
pinnedness in B.CF1. In particular no compatible structural 2-cells
are assumed in the definition of a pro-morphism. The following
argument supplies the coherence that is used, and then reconstructs
at a fixed source field. This avoids trying to make the transition
maps pinned.

\begin{claim}{B.PRO1}{Proved}{Compact coherence for ordinary pro-morphisms}
Let \(k\) have characteristic zero, let \(X,Y\) be finite-type
\(k\)-schemes, and let
\(\alpha:\mathcal X\to\mathcal Y\) be an ordinary pro-morphism
over \(\mathcal B\). There is a geometric morphism
\[
                 \alpha_k:X_{\mathrm{et}}\longrightarrow
                           Y_{\mathrm{et}}                 \tag{B.PRO.1}
\]
over \(BG_k\), compatible with the finite-level representatives of
\(\alpha\). Its underlying-space map is the map on inverse limits
defined by \(\alpha\), so it is pinned if \(\alpha\) is pinned.
Moreover finite constant base change, precomposition by a scheme
morphism, restriction to open subobjects, and slicing by a
quasi-compact separated étale object preserve the class of maps
obtained from ordinary pro-morphisms. For a slice defined at a
particular target level, the source model can be taken to be its
pullback by that level's chosen representative.
\end{claim}
\begin{proof}
Nilpotents do not affect the étale sites, so reduce the source.
At each target field \(F\), choose a representative
\(a_F:X_{E_F,\mathrm{et}}\to Y_{F,\mathrm{et}}\), where
\(E_F\supset F\) is finitely generated. Equality over the pro-base
permits one further source refinement making this a relative
morphism over \(BG_F\). Its structural 2-cell is unique:
\(G_F\) is slim by the opened Hoshi--Tsujimura Corollary~4.5,
and the image of each connected component's fundamental group in
\(G_F\) is open. For the latter assertion first map to \(G_{E_F}\);
the relative algebraic closure of \(F\) in the finitely generated
field \(E_F\) is finite. The same argument applies after every
finitely generated source refinement. Thus every absolute
comparison isomorphism at such a refinement is automatically
relative to \(BG_F\).

The opened \citep[Theorem~A.8 and Corollary~A.9]{CarlsonHaineWolf2024}
apply: the source is topologically noetherian and the target qcqs.
Every representative is coherent, and its inverse image takes
quasi-compact separated étale schemes to quasi-compact separated
étale schemes. Compose with \(X\to X_{E_F}\), obtaining functors
from the qcs étale sites of \(Y_F\) to that of \(X\).

Here is the finite-choice point that permits coherence. For two
qcs separated étale schemes \(U,V\) over a reduced noetherian
scheme \(X\) with finitely many irreducible components,
\(\operatorname{Isom}_X(U,V)\) is finite. An isomorphism is
determined on the finite fibres over the generic points: equality
there implies equality on \(U\), since \(U\) is reduced, those
fibres are dense, and \(V\to X\) is separated. The possible maps
of the finite generic étale algebras form a finite set.
Consequently the set of natural isomorphisms between two of the
preceding inverse-image functors is a closed subset of a product
of finite sets. It is compact, and the relative condition over
\(BG_F\) is closed as well.

For \(F\subset H\), ordinary pro-compatibility supplies an
isomorphism between the functor at \(F\) and the restriction of
the functor at \(H\), after a single finitely generated source
refinement and hence after base change to \(k\). The compact
space of relative comparison isomorphisms is therefore nonempty.
In each such compact space use the \emph{closure of the subset
of isomorphisms which descend after a finitely generated source
refinement}. This refinement is important: an arbitrary isomorphism
over \(k\) need not descend as an isomorphism of the entire
inverse-image functors. Membership in this closure says that on
every finite list of étale objects its components agree with those
of one whole comparison isomorphism descending at a finite stage.
Take the product of these closed spaces over all pairs of fields, and
impose naturality, identity, and the triple cocycle equations.
These are closed conditions. Every finite collection of them is
satisfiable: choose one common upper target field, compare its
representative with all the finitely many representatives after
one common finitely generated source refinement, and obtain all
comparisons by passing through this one representative. They
then satisfy every required finite cocycle equation. Compactness
gives a simultaneous coherent choice. This proves existence of
coherence; it has not been inserted into the hypothesis.

Every qcs étale object, every morphism between such objects, and
every finite diagram over \(Y\) descends to a finitely generated
field, by finite presentation. Finite covering families descend
too; the exact site-continuity statement is the opened
\citep[Lemma~59.51.2, Tag~0EYL, together with
Lemma~32.10.1, Tag~01ZM]{StacksProject}. We use the canonical
associative base-change identifications, or equivalently first
strictify this fixed diagram of sites. The coherently identified
inverse-image functors consequently
define a left exact covering-preserving functor on this site of
\(Y\), hence (B.PRO.1). The identical construction on finite étale
constant objects gives its structural isomorphism over \(BG_k\).
The comparison on open subobjects proves the assertion on spaces.

For clarity, slicing entails no additional coherence assumption.
Suppose \(V_F\to Y_F\) is qcs separated étale and put
\(U_{E_F}=a_F^*V_F\). For any higher target level, the preceding
comparisons identify the pullback of \(V\) with \(U\) over \(X\).
This is an isomorphism between finitely presented étale schemes,
so it descends after one finitely generated source refinement.
The resulting slice map represents a map of the model systems
\(\mathcal U\to\mathcal V\). For ordinary pro-compatibility,
approximate a selected comparison on the finite list consisting
of \(V\) and any specified arrows between slice objects by a
\emph{whole} comparison isomorphism descending at a finite stage.
This is exactly the closure property imposed above. Its component
on \(V\) agrees with the selected identification, so it induces
the required comparison of the entire slice geometric morphisms
after one refinement. Thus descending an individual isomorphism
of étale objects is not being substituted for descending a
comparison of functors. The same construction
over a finite étale arithmetic base gives finite constant base
change; cofinal fields have the form \(F'L_0\) for a fixed finite
field extension \(L_0/F_0\) defining that base change. Open
restriction and scheme precomposition are immediate. If the
original map is pinned, so are the slice maps: an étale
finite-type scheme has only closed points above a closed point.
\end{proof}

\begin{claim}{B.PRO2}{Proved}{Pro-points and a single-stage Kummer comparison}
For every characteristic-zero \(k\), pinned ordinary pro-morphisms
\(\mathcal B\to\mathcal Y\) over \(\mathcal B\) are precisely
the \(k\)-rational points of \(Y\). This statement is compatible
with finite extensions of \(k\).

In particular, suppose a fixed representative of a pinned
pro-morphism pulls the universal Kummer class back to
\(\xi\in H^1(X_E,\widehat\Z(1))\). If \(x\in X_E(L)\), with
\(L/E\) finite, and \(b\) is the value of its induced pro-point
in \(\mathbb G_m\), then there is \emph{one} finitely generated
field \(M\) containing \(L\) and \(b\) for which
\[
        \operatorname{res}_{M/L}(x^*\xi)=\kappa_M(b)
                \quad\hbox{in }\widehat{M^\times}.       \tag{B.PRO.2}
\]
The field \(M\) works simultaneously for all Kummer levels.
\end{claim}
\begin{proof}
Let the limit point of a pinned pro-point be \(y\in Y\).
It is closed, and a finite closed subscheme containing it descends
to some \(Y_{F_0}\). At every sufficiently large target field
\(F\), its image \(y_F\) is therefore closed, with
\(\kappa(y_F)/F\) finite. A representative
\(BG_E\to Y_{F,\mathrm{et}}\) has underlying image \(y_F\).
It factors through the closed subtopos
\(B G_{\kappa(y_F)}\): the inverse image of the complementary
open subobject is empty. After one source refinement it is
relative over \(BG_F\), with the unique structural 2-cell
established in B.PRO1. Since \(\kappa(y_F)/F\) is finite
separable, the finite étale slice description says that this
relative morphism is exactly an \(F\)-embedding
\(\kappa(y_F)\hookrightarrow E\). Thus it is an actual
\(E\)-valued point of \(Y_F\). Compatibility at a common
source refinement makes these points the same \(k\)-point;
the finite residue extension of the limit point consequently
has degree one. This also proves uniqueness of the pro-morphism.

For the second assertion base change the pro-morphism to the
finite extension \(kL/k\), using B.PRO1, and compose with the
point \(x\). The first assertion gives an actual value
\(b\in(kL)^\times\). Choose a finitely generated target field
\(H\subset kL\) containing \(b\) and all fields needed for the
displayed point. At target \(\mathbb G_{m,H}\), the pro-point's
representative has the closed underlying image \(T=b\).
After one refinement making it relative to \(BG_H\), it is
the constant-point morphism, as just proved. Its comparison
with the original representative at the original target is an
isomorphism of \emph{whole geometric morphisms} after one common
finitely generated source refinement \(M\). Pulling back all
the compatible universal Kummer torsors gives (B.PRO.2).
This proves simultaneous equality in the completed multiplicative
group, rather than a sequence of equalities over fields depending
on \(n\), or an equality only after restriction to \(G_{kL}\).
\end{proof}

\begin{claim}{B.PRO3}{Proved}{Intersection with a finitely generated field's Kummer completion}
Let \(A\subset B\) be finitely generated fields of characteristic
zero. If \(\xi\in\widehat{A^\times}\), \(b\in B^\times\), and
\[
                    \operatorname{res}_{B/A}\xi=\kappa_B(b),
                                                               \tag{B.PRO.3}
\]
then \(b\in A^\times\) and \(\xi=\kappa_A(b)\).
\end{claim}
\begin{proof}
Let \(A'\) be the relative algebraic closure of \(A\) in
\(B\); it is finite over \(A\). Take a normal projective
\(A'\)-model of the finitely generated regular extension
\(B/A'\). Every codimension-one valuation of this model is
trivial on \(A\). Applying its completed valuation to
(B.PRO.3) shows that the ordinary integer valuation of \(b\)
is zero. Normality then makes \(b\) a global invertible function
on this proper model. Such functions belong to its constant
field \(A'\), so \(b\in A'^\times\).

The map \(\widehat{A'^\times}\to\widehat{B^\times}\) is
injective: for every \(n\), an \(n\)-th root in \(B\) of an
element of \(A'\) is algebraic over \(A'\), hence already in
\(A'\). Thus (B.PRO.3) already holds over \(A'\). Taking the
norm to \(A\), with \(d=[A':A]\), gives
\[
                       d\xi=\kappa_A(N_{A'/A}b).          \tag{B.PRO.4}
\]

We record the elementary saturation argument, so no unsupported
finite-generation claim about \(A^\times\) is needed. In the
completion \(\widehat{A^\times}=\varprojlim_n
A^\times/(A^\times)^n\), if \(d\eta=0\), then its component
modulo \(n\) is represented by a \(d\)-th root of unity.
Indeed lift to level \(nd\): if \(a^d=c^{nd}\), then
\(a=c^n\zeta\) with \(\zeta^d=1\). The finite set
\(\mu_d(A)\) and compatibility over common multiples imply
\(\eta=\kappa_A(\zeta)\) for one \(\zeta\in\mu_d(A)\).
If \(d\xi=\kappa_A(a)\), reduction modulo \(d\) first gives
\(a=c^d\) with \(c\in A^\times\). Apply the preceding
torsion assertion to \(\xi-\kappa_A(c)\); hence \(\xi\)
is the Kummer class of an actual element \(c'\in A^\times\).
Apply this to (B.PRO.4). Finally Kummer maps of finitely generated
characteristic-zero fields are injective, by
\citep[Proposition~5.24 and Observation~6.9]{CarlsonHaineWolf2024}.
Equation (B.PRO.3) therefore makes \(c'=b\) in \(B\), proving
both assertions.
\end{proof}

\begin{claim}{B.PRO4}{Proved}{Reconstruction at one fixed source field}
Let \(\alpha\) be a pinned ordinary pro-morphism and choose a
relative representative
\[
                   a:X_{E,\mathrm{et}}\longrightarrow
                     Y_{F,\mathrm{et}}                 \tag{B.PRO.5}
\]
over \(BG_F\), with \(F\subset E\subset k\) finitely
generated. If \(X_E\) is normal, then (B.PRO.5) is induced by a
unique \(F\)-scheme morphism \(f_E:X_E\to Y_F\).
Its lift to \(Y_E\), after extension to \(k\), realizes the
point function of \(\alpha_k\). The same conclusion holds with
the absolutely weak normalization of \(X_E\) as source.
\end{claim}
\begin{proof}
\emph{1. Units over the fixed field.}
First take \(Y_F=\mathbb G_{m,F}\), and suppose \(X_E\) is
affine, geometrically integral, normal, and has an \(E\)-point.
The pullback \(\xi\) of the compatible Kummer classes lies in
\(\widehat{\mathcal O(X_E)^\times}\). For normal schemes this
follows from the Kummer sequence and finite generation of
\(\operatorname{Pic}(X_E)\), as stated in the opened
\citep[Proposition~5.22]{CarlsonHaineWolf2024}.
For every finite \(L/E\) and \(x\in X_E(L)\), B.PRO2 gives
one finitely generated \(M\supset L\) and a point value \(b\)
satisfying (B.PRO.2). B.PRO3 gives
\[
                         b\in L^\times,\qquad
                       x^*\xi=\kappa_L(b).              \tag{B.PRO.6}
\]
In particular the value at the chosen \(E\)-point is an actual
element of \(E^\times\); divide it out.

The normalized unit group is finitely generated free, by
\citep[Corollary~6.18]{CarlsonHaineWolf2024}. The evaluation of
the normalized completed unit at every algebraic point is an
actual algebraic unit by (B.PRO.6). The arithmetic argument of
Proposition~6.22 of that source now makes the completed unit
an actual unit of \(X_E\). Its hypotheses are precisely a
finitely generated characteristic-zero field (hence amenable
and admitting a nontrivial discrete valuation), the stated
finitely generated free unit group, and these actual point
evaluations. The proof works for normal sources as well:
Picard finite generation is Proposition~5.22, the unit assertion
is Corollary~6.18, and the nonconstant unit maps used in Claim~6.23
have image containing a nonempty open of \(\mathbb G_m\).
The source explicitly records this normal-source extension in
Remark~4.20. Equivalently, its full-rank integer valuation matrix
makes an integer multiple of each profinite exponent integral;
\(\widehat\Z\cap\Q=\Z\) then makes each exponent integral.
Thus \(\xi=\kappa(f)\) for \(f\in\mathcal O(X_E)^\times\).

This \(f\) realizes the value at \emph{every} point of
\(X(\overline k)\), not only the points algebraic over \(E\).
For such a point choose a finitely generated field of definition
\(L\supset E\) inside \(\overline k\). This extension need
not be finite. First base change the pro-system to the finite
extension \(kL/k\), then refine its source model from \(E\)
to \(L\). The point is now \(L\)-rational, so B.PRO2 applies
with the identity finite extension of the new source field
\(L\); B.PRO3 applies to the resulting finitely generated
comparison field. Its value \(b\) satisfies
\(\kappa_L(b)=\kappa_L(f(x))\); injectivity gives
\(b=f(x)\). The finite extension of \(k\) needed to contain
the point is permitted by B.PRO1.

Remove the rational-point and geometric-connectedness assumptions
by a finite Galois extension of \(E\) splitting the finitely
many geometric components and supplying a point on each. The
preceding proof gives units on these components. Their Kummer
classes are the pullback of the one fixed \(\xi\), so Galois
conjugates agree by Kummer injectivity. They descend to a unit
over \(E\). Passing to this finite extension is justified by
finite constant base change to its compositum with \(k\).
Finally a finite affine open cover handles a nonaffine normal
source; equality of the point functions glues the units uniquely.

\emph{2. Coordinates and the point function.}
For general \(Y_F\), inverse images of Zariski opens under
\(a\) are Zariski opens of \(X_E\), and the corresponding pro
restrictions are pinned. For a regular function \(h\) on an
affine target open, cover that open by \(D(h)\) and \(D(h-1)\).
Step~1 reconstructs the pulled-back units \(h\) and \(h-1\)
on their respective source opens over the same field \(E\).
On the overlap their difference is one, since this holds on
all \(\overline k\)-points and the source is reduced. They
therefore give a regular function on the entire inverse-image
open. The ring operations and all defining relations hold on
\(\overline k\)-points, hence as functions. The resulting maps
on affine target opens glue uniquely to
\(f_E:X_E\to Y_F\). Their point function is that of
\(\alpha_k\). The open inverse-image conditions hold as well,
because finite-type schemes over a field are Jacobson and
membership in an open is determined by the closed points of
its complement. This argument does not require the target to
be separated.

\emph{3. The entire functor at the same source level.}
For each qcs separated étale \(V_F\to Y_F\), the sheaf
\(a^*V_F\) is represented by a qcs separated étale scheme
\(U_E\to X_E\), by CHW A.8--A.9. It is normal. Slice the
pro-morphism, as in B.PRO1, using the fixed representative
\(a_V:U_{E,\mathrm{et}}\to V_{F,\mathrm{et}}\).
Step~2 gives a scheme map \(g_V:U_E\to V_F\) over this
same field \(E\), with the slice point function. This step
does not require the whole topos map of the slice to have been
reconstructed already. Composing to \(Y_F\) and comparing
point functions gives
\[
                 g_V\text{ over }Y_F
                    = f_E\circ(U_E\to X_E).
\]
It induces \(h_V:U_E\to X_E\times_{Y_F}V_F\).

The coherent realization (B.PRO.1) supplies the usual étale
slice fibre bijection: above every \(x\in X(\overline k)\),
the points of \(U\) are precisely the points of \(V\) above
\(\alpha_k(x)\). The map \(h_V\) realizes that bijection.
An étale morphism between étale schemes over \(X\) which is
bijective on every geometric closed fibre is an isomorphism:
the failure loci are constructible, so any nonempty such locus
over a field would contain a closed point. Thus \((h_V)_k\)
is an isomorphism, and faithful flatness makes \(h_V\) an
isomorphism over \(E\). This argument requires no additional
extension depending on \(V\).

For a morphism \(V_F\to W_F\), the two maps between the
reconstructed pullbacks agree on \(\overline k\)-points by
the natural slice functor of \(\alpha_k\); reducedness makes
them equal as scheme maps. The \(h_V\) consequently form a
natural isomorphism \(a^*\simeq f_E^*\) on the entire qcs
étale site. That site is a basis, so (B.PRO.5) is the induced
geometric morphism. This is equality at the fixed source field,
not just an equality of limit point functions or an objectwise
collection of equalities after unrelated extensions.

\emph{4. Faithfulness, automorphisms, and a seminormal source.}
We spell out the descent needed for the last assertion. For two
\(F\)-scheme maps from a reduced finite-type \(E\)-scheme to
\(Y_F\), an isomorphism of their relative geometric morphisms
implies equality of their pullbacks of all open subobjects and
all Kummer classes of units on target opens. Unit Kummer maps
on the source opens are injective: evaluate at closed points,
whose residue fields are finitely generated of characteristic
zero, and use their injectivity, followed by reducedness.
The coordinate argument of Step~2 therefore makes the two
scheme maps equal. This proves the required faithfulness even
when \(E/F\) is transcendental.

An actual such scheme map has no automorphisms as a geometric
morphism. At any source point \(x\), its residue field and
that of its image are finitely generated over \(\Q\); the
image of the source residue Galois group in the target residue
Galois group is open. The latter group is slim. Thus the
centralizer describing an automorphism on that stalk is trivial.
Geometric stalks are conservative, proving the assertion for
the whole inverse-image functor.

Replace \(X_E\) by its absolute weak normalization; in
characteristic zero this is finite-type seminormal. Its finite
normalization \(Z_E\to X_E\) is an h-cover. Apply the normal
case to \(Z_E\) and to a normalization of the reduced fibre
product \(Z_E\times_{X_E}Z_E\). The preceding faithfulness
makes the two maps agree after this normalization. Since that
normalization is an h-cover, h-separatedness in CHW Theorem~4.1
makes them agree on \((Z_E\times_{X_E}Z_E)^{\mathrm{awn}}\)
before applying the h-sheaf equalizer. The opened
\citep[Theorem~4.1]{CarlsonHaineWolf2024}, which states that
\(T\mapsto\operatorname{Hom}(T^{\mathrm{awn}},Y_E)\)
is an h-sheaf, descends them to a scheme map on \(X_E\).
Their topos isomorphisms are unique by the just-proved absence
of automorphisms and hence satisfy the cocycle condition.
Étale-topos h-descent, the opened Proposition~4.8 of the same
source, descends the isomorphism to (B.PRO.5) itself. This proves
the last assertion and uniqueness.
\end{proof}

\begin{claim}{B.PRO5}{Proved}{The ordinary pro-form for every characteristic-zero field}
For every field \(k\) of characteristic zero, the functor from
finite-type \(k\)-schemes localized at universal homeomorphisms
to their model pro-systems is fully faithful, with pinned
morphisms in the ordinary slice
\(\operatorname{Pro}(\mathsf T)_{/\mathcal B}\) as defined
above. In particular B9 is proved in arbitrary transcendence
degree; no coherent-relative morphism condition has been added
to its input.
\end{claim}
\begin{proof}
For a pinned ordinary pro-morphism \(\alpha\), choose a target
field \(F_0\) and a relative representative at a finitely
generated source field \(E\). B.PRO4 reconstructs it as a
localized scheme map \(X_E\to Y_{F_0}\). Pairing with the
structural map to \(E\) and extending to \(k\) gives
\(f:X^{\mathrm{awn}}\to Y\). At every other target field
\(F\), B.PRO4 reconstructs that level as well. Pro-compatibility
compares the two maps after one common finitely generated
source extension. The faithfulness proved in Step~4 of B.PRO4
makes their projections to \(Y_{F_0}\) equal; their structural
maps agree, so their lifts to \(Y_F\) and to \(Y\) agree.
At each target field the ordinary pro-class is therefore that
of \(f\). This is surjectivity on the Hom set in the ordinary
pro-category, rather than only on the corresponding limit
point functions.

Conversely, every localized finite-type \(k\)-scheme morphism
and its finite diagram of equations descend to a finitely
generated subfield, and hence define the model pro-morphism.
Its limit map preserves closed points. If two scheme morphisms
give the same pro-morphism, compare at a common target field,
refine once, and use the same fixed-field faithfulness; their
base changes to \(k\) are equal. Absolute weak normalizations
identify the localization, as in CHW Theorem~1.13 and
Proposition~2.28. This proves full faithfulness.
\end{proof}
